% GENERATED FILE. Edit canonical lessons, then run npm run generate:books.
% canonical-source-hash: fnv1a64:9bd7447a5cecaff8
% canonical-lessons: MR-C491-nidranas, MR-C491-ghorne, MR-C491-jambhai-dene, MR-C491-dulki-ghene, MR-C491-nijne

\chapter{Sleep and Naps}
\label{ch:mr-sleep-and-naps}
\hlchaptermodality{491}

\begin{chapteropening}
\textbf{By the end of this chapter:} \emph{I can talk about sleeplessness, snoring, yawning and napping.}

Complete the last lesson of chapter 491: I can talk about sleeplessness, snoring, yawning and napping.
\end{chapteropening}

\section[\texorpdfstring{\mr{निद्रानाश} (nidrānāś)}{निद्रानाश (nidrānāś)}]{\mr{निद्रानाश} (nidrānāś) — insomnia}

\label{lesson:MR-C491-nidranas}



\begin{quote}
Before the new one: say the Marathi for to apply heat, then the Marathi for to rub.
\end{quote}

\subsection*{You'll want to know: \mr{निद्रानाश}}
\textbf{\mr{निद्रानाश}} — \emph{nidrānāś} — \textquotedblleft{}insomnia\textquotedblright{}.

\begin{quote}
\mr{काळजीमुळे} \mr{त्याला} \mr{निद्रानाश} \mr{झाला}.
\end{quote}

\begin{quote}
\emph{kāḷjīmuḷe tyālā nidrānāś jhālā.}
\end{quote}

\begin{quote}
Worry gave him insomnia.
\end{quote}

Masculine.

\mr{निद्रा}, sleep, + \mr{नाश}, ruin.

\begin{cousinweb}
\textbf{\mr{निद्रानाश}} (\emph{nidrānāś}) is a learned compound built from two Sanskrit words, \emph{nidrā}, \textquotedblleft{}sleep\textquotedblright{} (CDIAL 7200), and \emph{nāśa}, \textquotedblleft{}loss, destruction\textquotedblright{} (Monier-Williams).
\end{cousinweb}

\begin{grammarlens}[title={What you've built}]
One more word for describing what you lived and felt.
\end{grammarlens}

\subsection*{Your turn}
\begin{itemize}
\raggedright
  \item \emph{Say it:} \emph{nidrānāś}
  \item \emph{Say it:} \emph{nidrānāś}, once more
  \item \emph{Say it:} \emph{kāḷjīmuḷe tyālā nidrānāś jhālā.} — Worry gave him insomnia.
  \item \emph{Recall:} say the Marathi for to apply heat, then the Marathi for to rub, then say \emph{nidrānāś} again
\end{itemize}

\begin{tcolorbox}[breakable,colback=teal!4,colframe=teal!35!black,title={Before you move on}]
What does \mr{निद्रानाश} mean? (\textquotedblleft{}Insomnia\textquotedblright{}.) Say it once more.
\end{tcolorbox}
\section[\texorpdfstring{\mr{घोरणे} (ghorṇe)}{घोरणे (ghorṇe)}]{\mr{घोरणे} (ghorṇe) — to snore}

\label{lesson:MR-C491-ghorne}



\begin{quote}
Before the new one: say the Marathi for to rub, then the Marathi for insomnia.
\end{quote}

\subsection*{You'll want to know: \mr{घोरणे}}
\textbf{\mr{घोरणे}} — \emph{ghorṇe} — \textquotedblleft{}to snore\textquotedblright{}.

\begin{quote}
\mr{आजोबा} \mr{रात्री} \mr{मोठ्याने} \mr{घोरतात}.
\end{quote}

\begin{quote}
\emph{ājobā rātrī moṭhyāne ghortāt.}
\end{quote}

\begin{quote}
Grandfather snores loudly at night.
\end{quote}

Past (he …): \mr{तो} \mr{घोरला}, \emph{to ghorlā}.

\mr{घोरण्याचा} \mr{आवाज}, the sound of snoring.

\begin{cousinweb}
\textbf{\mr{घोरणे}} (\emph{ghorṇe}) is a tadbhava: CDIAL lists Marathi \emph{ghorṇẽ}, \textquotedblleft{}to snore\textquotedblright{}, under Sanskrit \emph{ghurati}, \textquotedblleft{}cries out frighteningly\textquotedblright{} (CDIAL 4487).
\end{cousinweb}

\begin{grammarlens}[title={What you've built}]
One more word for describing what you lived and felt.
\end{grammarlens}

\subsection*{Your turn}
\begin{itemize}
\raggedright
  \item \emph{Say it:} \emph{ghorṇe}
  \item \emph{Say it:} \emph{ghorṇe}, once more
  \item \emph{Say it:} \emph{ājobā rātrī moṭhyāne ghortāt.} — Grandfather snores loudly at night.
  \item \emph{Recall:} say the Marathi for to rub, then the Marathi for insomnia, then say \emph{ghorṇe} again
\end{itemize}

\begin{tcolorbox}[breakable,colback=teal!4,colframe=teal!35!black,title={Before you move on}]
What does \mr{घोरणे} mean? (\textquotedblleft{}To snore\textquotedblright{}.) Say it once more.
\end{tcolorbox}
\section[\texorpdfstring{\mr{जांभई} \mr{देणे} (jāmbhaī deṇe)}{जांभई देणे (jāmbhaī deṇe)}]{\mr{जांभई} \mr{देणे} (jāmbhaī deṇe) — to yawn}

\label{lesson:MR-C491-jambhai-dene}



\begin{quote}
Before the new one: say the Marathi for insomnia, then the Marathi for to snore.
\end{quote}

\subsection*{You'll want to know: \mr{जांभई} \mr{देणे}}
\textbf{\mr{जांभई} \mr{देणे}} — \emph{jāmbhaī deṇe} — \textquotedblleft{}to yawn\textquotedblright{}.

\begin{quote}
\mr{कंटाळवाण्या} \mr{भाषणात} \mr{सगळे} \mr{जांभई} \mr{देत} \mr{होते}.
\end{quote}

\begin{quote}
\emph{kaṇṭāḷvāṇyā bhāṣṇāt sagḷe jāmbhaī det hote.}
\end{quote}

\begin{quote}
Everyone was yawning through the boring speech.
\end{quote}

Past, with -\mr{ने} (he …): \mr{त्याने} \mr{जांभई} \mr{दिली}, \emph{tyāne jāmbhaī dilī}.

\mr{जांभई}, a yawn; \mr{जांभई} \mr{येणे}, to feel a yawn coming, with the dative.

\begin{cousinweb}
\textbf{\mr{जांभई} \mr{देणे}} (\emph{jāmbhaī deṇe}) is a Marathi phrase built on \mr{जांभई}, \textquotedblleft{}a yawn\textquotedblright{}, which is inherited: CDIAL lists Old Marathi \emph{jāmbhaī}, \textquotedblleft{}yawning\textquotedblright{}, under \emph{*jṛmbhāyita} (CDIAL 5266); \mr{देणे}, \textquotedblleft{}to give\textquotedblright{}, is the light verb.
\end{cousinweb}

\begin{grammarlens}[title={What you've built}]
One more word for describing what you lived and felt.
\end{grammarlens}

\subsection*{Your turn}
\begin{itemize}
\raggedright
  \item \emph{Say it:} \emph{jāmbhaī deṇe}
  \item \emph{Say it:} \emph{jāmbhaī deṇe}, once more
  \item \emph{Say it:} \emph{kaṇṭāḷvāṇyā bhāṣṇāt sagḷe jāmbhaī det hote.} — Everyone was yawning through the boring speech.
  \item \emph{Recall:} say the Marathi for insomnia, then the Marathi for to snore, then say \emph{jāmbhaī deṇe} again
\end{itemize}

\begin{tcolorbox}[breakable,colback=teal!4,colframe=teal!35!black,title={Before you move on}]
What does \mr{जांभई} \mr{देणे} mean? (\textquotedblleft{}To yawn\textquotedblright{}.) Say it once more.
\end{tcolorbox}
\section[\texorpdfstring{\mr{डुलकी} \mr{घेणे} (ḍulkī gheṇe)}{डुलकी घेणे (ḍulkī gheṇe)}]{\mr{डुलकी} \mr{घेणे} (ḍulkī gheṇe) — to take a nap}

\label{lesson:MR-C491-dulki-ghene}



\begin{quote}
Before the new one: say the Marathi for to snore, then the Marathi for to yawn.
\end{quote}

\subsection*{You'll want to know: \mr{डुलकी} \mr{घेणे}}
\textbf{\mr{डुलकी} \mr{घेणे}} — \emph{ḍulkī gheṇe} — \textquotedblleft{}to take a nap\textquotedblright{}.

\begin{quote}
\mr{जेवणानंतर} \mr{आजोबा} \mr{थोडी} \mr{डुलकी} \mr{घेतात}.
\end{quote}

\begin{quote}
\emph{jevṇānantar ājobā thoḍī ḍulkī ghetāt.}
\end{quote}

\begin{quote}
After lunch Grandfather takes a short nap.
\end{quote}

Past, with -\mr{ने} (he …): \mr{त्याने} \mr{डुलकी} \mr{घेतली}, \emph{tyāne ḍulkī ghetlī}.

\mr{डुलकी}, a short doze; \mr{डुलकी} \mr{लागणे}, to nod off.

\begin{cousinweb}
\textbf{\mr{डुलकी} \mr{घेणे}} (\emph{ḍulkī gheṇe}) is a Marathi phrase built on \mr{डुलकी}, \textquotedblleft{}a doze\textquotedblright{}, whose origin is not settled, with \mr{घेणे}, \textquotedblleft{}to take\textquotedblright{}.
\end{cousinweb}

\begin{grammarlens}[title={What you've built}]
One more word for describing what you lived and felt.
\end{grammarlens}

\subsection*{Your turn}
\begin{itemize}
\raggedright
  \item \emph{Say it:} \emph{ḍulkī gheṇe}
  \item \emph{Say it:} \emph{ḍulkī gheṇe}, once more
  \item \emph{Say it:} \emph{jevṇānantar ājobā thoḍī ḍulkī ghetāt.} — After lunch Grandfather takes a short nap.
  \item \emph{Recall:} say the Marathi for to snore, then the Marathi for to yawn, then say \emph{ḍulkī gheṇe} again
\end{itemize}

\begin{tcolorbox}[breakable,colback=teal!4,colframe=teal!35!black,title={Before you move on}]
What does \mr{डुलकी} \mr{घेणे} mean? (\textquotedblleft{}To take a nap\textquotedblright{}.) Say it once more.
\end{tcolorbox}
\section[\texorpdfstring{\mr{निजणे} (nijṇe)}{निजणे (nijṇe)}]{\mr{निजणे} (nijṇe) — to lie down; to fall asleep}

\label{lesson:MR-C491-nijne}



\begin{quote}
Before the new one: say the Marathi for to yawn, then the Marathi for to take a nap.
\end{quote}

\subsection*{You'll want to know: \mr{निजणे}}
\textbf{\mr{निजणे}} — \emph{nijṇe} — \textquotedblleft{}to lie down; to fall asleep\textquotedblright{}.

\begin{quote}
\mr{बाळ} \mr{आईच्या} \mr{कुशीत} \mr{निजलं}.
\end{quote}

\begin{quote}
\emph{bāḷ āīcyā kuśīt nijlaṁ.}
\end{quote}

\begin{quote}
The baby fell asleep in its mother's arms.
\end{quote}

Past (he …): \mr{तो} \mr{निजला}, \emph{to nijlā}.

A softer, older word beside \mr{झोपणे}.

\begin{cousinweb}
\textbf{\mr{निजणे}} (\emph{nijṇe}) is a tadbhava, inherited from Sanskrit \emph{nidrāyati}, \textquotedblleft{}falls asleep\textquotedblright{}; CDIAL lists Marathi \emph{nijṇẽ} (CDIAL 7201).
\end{cousinweb}

\begin{grammarlens}[title={What you've built}]
One more word for describing what you lived and felt.
\end{grammarlens}

\subsection*{Your turn}
\begin{itemize}
\raggedright
  \item \emph{Say it:} \emph{nijṇe}
  \item \emph{Say it:} \emph{nijṇe}, once more
  \item \emph{Say it:} \emph{bāḷ āīcyā kuśīt nijlaṁ.} — The baby fell asleep in its mother's arms.
  \item \emph{Recall:} say the Marathi for to yawn, then the Marathi for to take a nap, then say \emph{nijṇe} again
\end{itemize}

\begin{tcolorbox}[breakable,colback=teal!4,colframe=teal!35!black,title={Before you move on}]
What does \mr{निजणे} mean? (\textquotedblleft{}To lie down; to fall asleep\textquotedblright{}.) Say it once more.
\end{tcolorbox}
